% TOP_PROBLEM: {"id": 30001926, "title": "Quantum Unique Ergodicity on Negatively Curved Manifolds", "category_id": 16, "category": {"id": 16, "name": "physics", "display_name": "Mathematical Physics", "description": "Problems at the intersection of mathematics and physics.", "slug": "physics", "order_index": 16, "created_at": "2026-07-31T15:26:25.671Z"}, "status": "partially_solved"}
% FIRST_POSTED: 2026-09-27
% !TeX program = lualatex
% Compile twice: lualatex 30001926.tex
% All references and prose are embedded; no BibTeX database or external inputs.
\documentclass[11pt,a4paper]{article}
\usepackage[margin=24mm,headheight=14pt]{geometry}
\usepackage{fontspec}
\setmainfont{Latin Modern Roman}
\setsansfont{Latin Modern Sans}
\setmonofont{Latin Modern Mono}
\usepackage{amsmath,amssymb,mathtools}
\usepackage{unicode-math}
\setmathfont{Latin Modern Math}
\usepackage{microtype}
\usepackage{enumitem,xurl,needspace}
\usepackage[dvipsnames]{xcolor}
\usepackage{hyperref}
\hypersetup{unicode=true,colorlinks=true,linkcolor=MidnightBlue,urlcolor=MidnightBlue,
 pdftitle={Research attempt: Problem 30001926},pdfauthor={Research notebook},bookmarksnumbered=true}
\usepackage{bookmark}
\usepackage{tocloft}
\setlength{\cftsecnumwidth}{3em}
\cftsetpnumwidth{2.5em}
\cftsetrmarg{4em}
\usepackage{titlesec}
\titleformat{\section}{\Large\bfseries\raggedright}{\thesection}{0.7em}{}
\usepackage{fancyhdr}
\pagestyle{fancy}\fancyhf{}
\fancyhead[L]{\small\sffamily Open Problems Math | Research notebook}
\fancyhead[R]{\small Problem 30001926 | 27 September 2026}
\fancyfoot[C]{\thepage}
\setlength{\parindent}{0pt}
\setlength{\parskip}{3pt plus 1pt}
\setlength{\emergencystretch}{3em}
\setcounter{tocdepth}{1}
\setcounter{secnumdepth}{1}
\setlist[itemize]{leftmargin=3em,itemsep=5pt,topsep=4pt,parsep=0pt}
\allowdisplaybreaks
\newcommand{\N}{\mathbb{N}}
\newcommand{\Z}{\mathbb{Z}}
\newcommand{\Q}{\mathbb{Q}}
\newcommand{\R}{\mathbb{R}}
\newcommand{\C}{\mathbb{C}}
\newcommand{\F}{\mathbb{F}}
\newcommand{\A}{\mathbb{A}}
\newcommand{\PP}{\mathbb P}
\newcommand{\E}{\mathbb{E}}
\newcommand{\eps}{\varepsilon}
\newcommand{\dd}{\,\mathrm d}
\newcommand{\sref}[2]{\hyperlink{source:#1:#2}{[S#2]}}
\newcommand{\eref}[2]{\hyperlink{extra:#1:#2}{[E#2]}}
% Turn the next switch off for a shorter reading copy. The quotations preserve
% every exactTarget field, including qualifications omitted from a short summary.
\newif\ifshowcatalogue
\showcataloguetrue
\newcommand{\cataloguescope}[1]{\ifshowcatalogue
 \paragraph{Exact scope in the supplied catalogue.}
 \begin{quote}\small #1\end{quote}\fi}
\usepackage{etoolbox}
\AtBeginEnvironment{itemize}{\small}
\numberwithin{equation}{section}
% Standalone export. Original research content preserved.
\begin{document}
\setcounter{section}{0}
% ================================================================
% problemId: 30001926
% Canonical problem ID: 30001926
% exactTarget SHA-256: b51adc4452620b773144ee6c07e676f1d1399f861adb58011c072e13ed28f31e
\clearpage
\section*{\texorpdfstring{Quantum Unique Ergodicity on Negatively Curved Manifolds}{Quantum Unique Ergodicity on Negatively Curved Manifolds}}\label{30001926}
\subsection{Definitions and mathematical statement}
Let $-\Delta_g u_j=\lambda_j^2u_j$, with $\|u_j\|_{L^2(M)}=1$ and $\lambda_j\to\infty$. Semiclassical measures are weak limits of the distributions
\[
 a\longmapsto\langle\operatorname{Op}_{1/\lambda_j}(a)u_j,u_j\rangle
\]
on the unit cosphere bundle $S^*M$. Liouville measure is its normalized canonical symplectic-volume measure. Quantum unique ergodicity requires every such limit to be Liouville for every compact negatively curved $M$. In particular, $|u_j|^2\,d\operatorname{vol}_g$ must tend to normalized volume. A density-one subsequence result does not exclude exceptional concentrating sequences.
\par\smallskip\textit{Formulation and concept references:} \sref{30001926}{1}.
\cataloguescope{Let (M,g) be a compact negatively curved Riemannian manifold. Prove that every complete high-energy sequence of L2-normalized Laplace eigenfunctions has semiclassical measures converging to Liouville measure on the unit cosphere bundle; consequently, the position densities converge weakly to normalized Riemannian volume.}
\subsection{Short English statement}
Must every high-energy eigenfunction sequence on a compact negatively curved manifold become uniformly distributed in phase space?
% BEGIN RESEARCH ADDITION 112
\subsection{Research attempt}

\paragraph{Current frontier and source audit.}
The cited talk is dated 5 January 2026. Its slide numbered 6 is PDF page 16
because the PDF includes overlays. It distinguishes quantum ergodicity,
QUE, and entropy bounds. Its later slides restrict full-support theorems
to surfaces and discuss separate higher-dimensional locally symmetric
results. The slide's abbreviated arithmetic statement must retain the
additional arithmetic eigenfunction hypotheses, which are absent here.
None of these statements identifies every quantum limit with Liouville
measure on every negatively curved manifold. \sref{30001926}{1}

Dyatlov--Jin--Nonnenmacher prove full support on compact surfaces with
Anosov geodesic flow, including variable negative curvature. Their
Theorem~2 is a quantitative microlocal control estimate. Kim--Miller's
2025 manuscript treats real hyperbolic manifolds in higher dimensions,
obtaining a support restriction involving an immersed totally geodesic
submanifold. These results do not supply the uniform equidistribution
estimate pursued below. \eref{30001926}{1}\eref{30001926}{2}

\paragraph{Attack route.}
The issue is an exceptional sequence, not an average eigenfunction. We
therefore replace averaged quantum variance by the norm of the observable
compressed to each exact eigenspace. Three possible upgrades are examined:
high Schatten moments; time-averaging followed by a trace estimate; and
combining entropy with full support. The first yields a precise sufficient
estimate. The second exposes a logarithmic-time obstruction. The third
fails even at the level of invariant probability measures.

\paragraph{Attempt: exact spectral compression.}
Fix a self-adjoint order-zero pseudodifferential operator $A$ whose
homogeneous principal symbol $a$ has Liouville mean zero on $S^*M$.
Write $P=\sqrt{-\Delta_g}$, and let $\Pi_\lambda$ project onto the complete
eigenspace of $P$ with eigenvalue $\lambda$. For $R>0$, set
\[
 \Pi_R=\sum_{R\leq\lambda<2R}\Pi_\lambda,\quad
 N_R=\operatorname{rank}\Pi_R,\quad
 D_R(A)=\sum_{R\leq\lambda<2R}\Pi_\lambda A\Pi_\lambda.
\]
All operators in the last display are now finite-dimensional on
$\operatorname{ran}\Pi_R$. The supremum over \emph{all} normalized eigenvectors
in this window is exactly
\begin{equation}\label{eq:30001926:block}
 \sup_{\substack{Pu=\lambda u,\ \|u\|=1\\ R\leq\lambda<2R}}
 |\langle Au,u\rangle|=\|D_R(A)\|.
\end{equation}
Indeed each block is self-adjoint and its extreme Rayleigh quotient is an
eigenvalue of that block. Such an extremizing vector is still an exact
Laplace eigenfunction, even if the Laplace eigenvalue is multiple.

Thus QUE for every sequence is equivalent to $\|D_R(A)\|\to0$ for every
such $A$. For sufficiency it is enough to test a countable symbol-dense
collection: uniform order-zero norm bounds extend the conclusion to all
smooth test symbols, lower-order errors vanish at high energy, and
compactness of semiclassical measures then gives the unique limit.
For necessity, failure of the norm convergence selects genuine
eigenfunctions violating QUE through \eqref{eq:30001926:block}. This formulation
avoids the unjustified choice of a preferred basis in a multiple eigenspace.

\textbf{High-moment sufficient lemma.}
Suppose for a given $A$ there are integers $k_R\geq1$ and numbers
$\epsilon_R\to0$ such that
\begin{equation}\label{eq:30001926:moment}
 \operatorname{Tr}|D_R(A)|^{2k_R}\leq N_R\epsilon_R^{2k_R},
 \qquad \sup_R\frac{\log\max(2,N_R)}{k_R}<\infty.
\end{equation}
Then $\|D_R(A)\|\to0$. This is immediate from
\[
 \|D_R(A)\|\leq
 (\operatorname{Tr}|D_R(A)|^{2k_R})^{1/(2k_R)}
 \leq N_R^{1/(2k_R)}\epsilon_R.
\]
No interchange of a growing moment and a spectral limit is hidden here:
\eqref{eq:30001926:moment} is the explicit unproved quantitative hypothesis.
If it holds for every $A$ above, the catalogue target follows in every
dimension to which that hypothesis applies.

A fixed-moment estimate is inadequate. The matrices
$D_N=\operatorname{diag}(1,0,\ldots,0)$ satisfy
$N^{-1}\operatorname{Tr}|D_N|^{2k}=N^{-1}\to0$ for every fixed $k$, but
$\|D_N\|=1$. This is an exact model of a single exceptional state missed
by all fixed moments. It is not asserted to arise from a Laplacian.

\paragraph{Attempt: time-averaging and the scale of the missing theorem.}
Define the bounded self-adjoint operator
\[
 A_T=\frac1T\int_0^T e^{itP}Ae^{-itP}\,dt.
\]
Since scalar phases cancel on each exact eigenspace,
$\Pi_\lambda A_T\Pi_\lambda=\Pi_\lambda A\Pi_\lambda$ for every $T>0$.
Let $B_{R,T}=\Pi_R A_T\Pi_R$. Block pinching is an average of finitely
many unitary conjugations: attach distinct roots of unity to the different
eigenvalue blocks and average powers of the resulting diagonal unitary.
Convexity and unitary invariance of each Schatten norm consequently imply
\begin{equation}\label{eq:30001926:pinching}
 \operatorname{Tr}|D_R(A)|^{2k}
 \leq\operatorname{Tr}|B_{R,T}|^{2k}\qquad(k\geq1).
\end{equation}
This step is exact, independent of Egorov's theorem, multiplicities, and
spectral spacings.

For comparison, suppose a genuinely uniform quantum estimate of the form
\begin{equation}\label{eq:30001926:desired}
 \operatorname{Tr}|B_{R,T_R}|^{2k_R}
 \leq N_R\left(\frac{Ck_R}{T_R}\right)^{k_R},\qquad
 k_R=\lceil\log\max(2,N_R)\rceil,
 \qquad \frac{T_R}{\log R}\longrightarrow\infty
\end{equation}
were available. The spectral counting bound $N_R=O(R^{\dim M})$ makes
$k_R=O(\log R)$; hence \eqref{eq:30001926:desired} implies
\eqref{eq:30001926:moment} with $\epsilon_R=(Ck_R/T_R)^{1/2}\to0$.
The required conclusion would follow. Any remainder must likewise have
$2k_R$-th root tending to zero after removing the factor $N_R$; a bare
$o(N_R)$ remainder is not sufficient for these growing moments.

The attempted derivation is to replace $A_T$ microlocally by the classical
average $T^{-1}\int_0^T a\circ\varphi^t\,dt$, estimate its high moments,
and use a trace formula. But derivatives of $a\circ\varphi^t$ grow
exponentially. The available hyperbolic semiclassical methods in the cited
control theorem work with carefully controlled logarithmic propagation
times, not the superlogarithmic uniform high-moment estimate
\eqref{eq:30001926:desired}. Even a favorable Gaussian moment ansatz
$(Ck/T)^k$ at $T=c\log R$, $k\asymp\log N_R$, gives only a constant bound
on $\epsilon_R$, not convergence to zero. Neither changing the order of
limits nor ignoring symbol seminorms repairs that calculation.
\eref{30001926}{1}

\paragraph{Stress test: full support plus entropy still leaves room.}
On a curvature $-1$ compact surface, let $\mu_\gamma$ be the invariant
probability measure on a periodic geodesic (averaged over both orientations
when imposing time reversal). For $0<\theta\leq1/2$, set
\[
 \mu_\theta=(1-\theta)\mu_L+\theta\mu_\gamma.
\]
This measure has full support, is invariant, and is not Liouville. Entropy
is affine on invariant measures, so
$h_{\mathrm{KS}}(\mu_\theta)=1-\theta\geq1/2$.
It therefore passes the full-support and half-entropy necessary conditions
recalled in the source. \sref{30001926}{1}
We have \emph{not} constructed eigenfunctions realizing $\mu_\theta$.
The point is that those two necessary conditions, considered by themselves,
cannot force the desired conclusion.

Similarly, classical ergodicity makes a zero-mean time average small in
$L^2(\mu_L)$, not uniformly at every phase point. If $a$ has nonzero average
on a periodic orbit, its time average there remains nonzero. Replacing a
trace norm by an operator norm solely on the strength of classical
ergodicity is therefore exactly the unproved step, not a harmless estimate.

\paragraph{Outcome.}
\textbf{Outcome: Conditional reduction.}

The notebook proves the basis-independent block criterion, the pinching
inequality, and a growing-moment sufficient condition. It also identifies
the time/moment scale that a direct mixing-plus-trace strategy would have
to reach and supplies exact countertests to its tempting weaker versions.
The needed quantum high-moment estimate is unproved. No new QUE theorem,
exceptional Laplace eigenfunction, or counterexample manifold is claimed.
% END RESEARCH ADDITION 112
\subsection{Sources}
\begin{itemize}\raggedright
\item[\textnormal{[S1]}]\hypertarget{source:30001926:1}{} Dyatlov, Control of Eigenfunctions on Negatively Curved Manifolds, QUE and entropy bounds, slide 6 (2026).
\url{https://math.mit.edu/~dyatlov/files/2026/chyptalk.pdf}.
{\small\textit{Catalogue formulation source.}}
% BEGIN NEW REFERENCES 112
\item[\textnormal{[E1]}]\hypertarget{extra:30001926:1}{} Semyon Dyatlov, Long Jin and
St\'ephane Nonnenmacher, \emph{Control of eigenfunctions on surfaces of variable
curvature}, arXiv:1906.08923v2; J. Amer. Math. Soc. 35 (2022), 361--465.
Theorems 1--2, Sections 3--5 and Appendix A.3.
\url{https://arxiv.org/html/1906.08923v2}.
\item[\textnormal{[E2]}]\hypertarget{extra:30001926:2}{} Elena Kim and Nicholas Miller,
\emph{Semiclassical Measures on Hyperbolic Manifolds}, arXiv:2503.01528v1
(2025), introduction and main support theorem.
\url{https://arxiv.org/html/2503.01528v1}.
% END NEW REFERENCES 112
\end{itemize}

\end{document}
